% sections/intro.tex -- Section 1 (paper_v2, cut in C3 from the full version's introduction).
% B1 (25 November) brought in the author's review of the Vietnamese introduction (paper_v2_vn,
% L10): the three reasons as sentences, the question against manual design, two ways with random
% draws as the control, a verdict that depends on the context, and no sentence on cleaning the
% data, since the first contribution gives it. The claims stay within the full version's, and every
% number is \val{} from numbers.tex.
\section{Introduction}
\label{sec:intro}

Recommender systems often know more than who bought what: products belong to categories,
customers buy some of them together, and users review them. A heterogeneous graph holds all of
these relations, and meta-paths are a common way to read it: \code{PUP} (product--user--product)
links two products reviewed by the same user. HAN \citep{wang2019heterogeneous} turns each
meta-path into a graph and weighs the paths with attention, and many recommenders are built on
meta-paths \citep{hu2018leveraging,shi2020hgrec,mu2021fhrec,yan2021attention}. The choice of paths
matters: the performance of meta-path-based recommenders is ``highly sensitive to the meta-paths
they use'' \citep{ning2022automatic}.

The paths are usually chosen by hand, which takes expert knowledge and does not carry over from
one dataset to the next \citep[][App.~D]{han2020genetic,li2024longrange}, and automatic discovery
aims to take this choice over (Section~\ref{sec:related}). LMSPS \citep{li2024longrange}, a one-shot
search over a large space of candidate paths, reports that its searched paths raise HAN's Micro-F1
by \val{lmsps-t4-gain-min} to \val{lmsps-t4-gain-max} points on four node-classification datasets
\citep[Table~4]{li2024longrange}. To our knowledge, whether they also help HAN as a recommender has
not been tested, and the answer is not obvious for three reasons. First, LMSPS judges paths with its
own search network rather than with HAN, so the best paths for that network need not be the best for
HAN. Second, studies of neural architecture search found that searches often do no better than
random choices \citep{li2019random,yu2020evaluating}. Third, the search costs GPU hours that a
hand design does not need.

We therefore ask: \emph{compared with manual design, is meta-path search worth its cost for
HAN-based product recommendation?} On Amazon-Meta \citep{leskovec2007viral}, we compare two ways of
choosing HAN's meta-paths from the same \val{search-cands} candidates. The first is by hand, with
sets designed with AI assistance from stated principles and fixed before any result. The second is
search by LMSPS, adapted to user--product link prediction and run \val{search-runs} times per set
size, since one run cannot rank its best sets. To see how much the choice of paths matters, we use
random draws from the same candidates as the control. Every set is trained in HAN and metapath2vec
\citep{dong2017metapath2vec} and tested on the same held-out edges. The main metric is sampled
Hit@10 (fixed in advance), with full ranking as a check. Popularity and a type-blind GAT serve as
baselines.

On this schema the answer depends on the context, since the searched sets come fairly close to the
manual ones. On the metric fixed in advance, they trail the manual sets by \val{han-4-hit-relabs}\%
and \val{han-6-hit-relabs}\% in HAN at four and six paths (\val{han-4-hit-pt} and
\val{han-6-hit-pt}) and by \val{mtv-4-hit-relabs}\% and \val{mtv-6-hit-relabs}\% in metapath2vec. In
HAN they tie on full ranking, after about \val{search-total-h} GPU hours of search. Manual design is
therefore the better choice when an expert can design the paths in less time than the search takes,
while search gives nearly equivalent sets when no expert is at hand. The searched sets lose mostly
on books and music, the products their category paths cannot reach. The random control shows that
the choice of paths matters little on this schema: in HAN, random draws do at least as well as both
ways. GAT also beats both meta-path models. Two limits remain. First, with a shorter path-weight
warm-up in HAN, the searched set beats the manual one at six paths in a check with a single seed
(Section~\ref{sec:res-sensitivity}). Second, this small schema leaves untested the case that
automatic search is meant for: schemas too large to design by hand or to sample at random.

Our contributions are:
\begin{itemize}
  \item A cleaned Amazon-Meta for link prediction, in which merging duplicate editions cuts the
        test edges with a copy in training from \val{before-testleak-pct}\% to
        \val{after-testleak-pct}\% (Section~\ref{sec:data}).
  \item LMSPS adapted to link prediction with leakage controls, and a way to combine many search
        runs into one decision (Section~\ref{sec:method}).
  \item A controlled comparison of searched and manual path sets in HAN and metapath2vec, with
        random draws as the control and with the sets and the main metric fixed in advance
        (Sections~\ref{sec:setup} and~\ref{sec:results}).
  \item A cost--performance view of search against manual design, with path-level evidence for
        where the searched sets lose (Section~\ref{sec:results}).
\end{itemize}
